在任何单个 \(L\)-自变量中，分离变量只有有限多个。给每个变量分配 \(T_1/2\) 这一单项配额中的一个固定份额；不要在相继的移位处重新分配该配额。行估计中所用的剖面窗口与有限高度阶都在 \(N\) 之前固定。引理~\ref{lem:smooth-calculus} 在增大 \(N\) 时只对未扭变的分离剖面求导，这便证明了最后的断言。
\end{proof}

\subsection{同时的饱和见证}

在一个选定的零点处，下面的截断逆构造产生一个乘积，其平方大小至少为 $U^{\delta(r+m)-\epsilon}$。
逐点估计将同一乘积界为 $U^{\delta(r+\min(m,1-m))+\epsilon}$。由于 $\delta\ge1/50$，
这两个界迫使 $m\le1/2+O(\epsilon)$。于是每个因子都在指定的损失范围内达到其自身的逐点指数；
这就是下一命题中所断言的饱和性。

\begin{proposition}[两个饱和见证]\label{prop:detector-witness}
设 \(u\) 有一个分箱 \((i,a)\)，其中 \(a>51/100\)，并假设引理~\ref{lem:detector-dyads} 的损失与高度假设。
对每个 \(t\in[1,3/2]\)，以及每个指定的 \(\epsilon>0\)，
在缩减其初步损失之后，存在
\(\psi\in\mathcal X_u\)、如引理~\ref{lem:buffered-bins} 那样的一个零点 \(\rho=\sigma+i\gamma\)、
二进长度 \(D=U^r\)、\(N=U^m\)，以及公共行特征 \(\psi\) 的两个多项式 \(M_r,S_m\)，使得
\begin{equation}\label{eq:detector-note}
 r\le t+O(1/\log U),\qquad
 r+m\ge t-O(1/\log U),\qquad
 |M_rS_m|^2\gg_{\mathcal A,\epsilon}
 U^{\delta(r+m)-\epsilon}.
\end{equation}
它们的长度与各自的值也满足
\begin{equation}\label{eq:saturated-witness}
 \begin{gathered}
 t-\tfrac12-O(\epsilon)\le r\le t+O(1/\log U),
 \qquad 0\le m\le\tfrac12+O(\epsilon),\\
 |M_r|^2\gg_{\mathcal A,\epsilon}U^{\delta r-\epsilon},
 \qquad
 |S_m|^2\gg_{\mathcal A,\epsilon}U^{\delta m-\epsilon}.
 \end{gathered}
\end{equation}
在所述参数范围内，\(O(\epsilon)\) 项中的常数是绝对的。两个多项式具有相同的扭变高度
\(\gamma-\nu\)，其中 \(|\nu|\le cT_1\)，\(c>0\) 固定且取自累积频率配额之内。
它们的未扭变剖面构成一个一致光滑的环形族。
对每一行，两个见证都使用一个表示 \(\psi_{u,\vartheta,\varsigma}\)，由一个固定的有限集合
\((\vartheta,\varsigma)\)，取自固定有限集合
\(\Theta\times\{-1,1\}\) 中的标号选取。表示本身随 \(u\) 变化。
二进对同样可以逐行地从 \(O((\log U)^2)\) 种可能性中分别选取。
\end{proposition}

\begin{proof}
截断逆与 Gamma 积分构造是经典零点探测器的一种形式；比较 \cite[附录 C]{MP}
与 \cite[第~13.1 节]{GM}。此处所需的带缓冲的行相关矩形与
同时下界将在下面建立。选取由该分箱给出的零点 \(\rho\)，并令
\[
 D_*=U^t,\qquad Y_*=U^{20}.
\]
设 \(V_{\le}\) 为一个固定的光滑函数，在
\([0,1]\) 上等于一、在 \([2,\infty)\) 上等于零。定义
\[
 C_\psi(s)=\sum_l\mu(l)\psi(l)V_{\le}(q_l/D_*)q_l^{-s},
 \qquad
 J=\frac1{2\pi i}\int_{(2)}
       Y_*^z\Gamma(z)L_{\rm orig}(\rho+z,\psi)
       C_\psi(\rho+z)\,dz.
\]
行特征是非主特征，故 \(L_{\rm orig}\) 为整函数。
将直线移至 \(\Re z=-1/4\) 只穿过
\(\Gamma(z)\) 在零处的极点，且其留数为零，因为
\(L_{\rm orig}(\rho,\psi)=0\)。方程~\eqref{eq:hecke-growth} 中的全局带形界，
连同被删除的 Euler 因子，将新直线界为
\[
 U^{-5}U^2D_*^{\,1-\sigma+1/4+o(1)}T_1^2.
\]
事实上 \(\Re(\rho-1/4)\ge26/100>0\)，故被删除的乘积具有
任意小的 \(U\)-幂；刻意取弱的
\(U^2\) 包含了本原导子界。有限
多项式利用理想计数由所示 \(D_*\) 的幂界出。最后，所加高度中的固定多项式
对 \(\Gamma(-1/4+iv)\) 可积，因为
\((1+|\gamma+v|)^2\le(1+|\gamma|)^2(1+|v|)^2\)。
对 \(1\le t\le3/2\) 一致地有
\[
 -3+t(5/4-\sigma)
 \le-3+\frac32\left(\frac54-\frac{51}{100}\right)
 =-\frac{189}{100}.
\]
于是 \(J=o(1)\)。这一 Gamma 积分只使用全局上
界，不消耗任何倒数高度配额。

在原始直线上，绝对收敛性与指数的 Mellin 公式
给出
\[
 J=\sum_{l,k}\mu(l)\psi(lk)V_{\le}(q_l/D_*)
       (q_lq_k)^{-\rho}e^{-q_lq_k/Y_*}.
\]
对 \(1<q_n\le D_*\)，\(\psi(n)q_n^{-\rho}\) 的系数
为 \(\sum_{l\mid n}\mu(l)=0\)。对 \(n=1\) 它为一。
当 \(q_n>D_*\) 时，\(V_{\le}(2q_n/D_*)=0\)。
因此插入 \(1-V_{\le}(2q_lq_k/D_*)\) 恰好
移除单位贡献 \(e^{-1/Y_*}=1+o(1)\)。
所得尾部具有绝对值 \(1+o(1)\)。
用这样一个固定的光滑终端截断函数乘它——该函数对
\(q_lq_k\le U^{21}\) 等于一、对 \(q_lq_k\ge2U^{21}\) 等于零——将其改变
\(o(1)\)；这由 \(Y_*=U^{20}\)、指数衰减
以及理想计数可得。

用固定的光滑二进划分将 \(q_l\asymp D\)、\(q_k\asymp N\)
划分。相关的对共有 \(O((\log U)^2)\) 个，
其支集满足
\[
 D\ll D_*,\qquad DN\gg D_*,\qquad DN\ll U^{21}.
\]
以 \(x=q_l/D\)、\(y=q_k/N\)，一对的剖面为
\[
 W_1(x)V_{\le}(Dx/D_*)\,W_2(y)\,H_{D,N}(xy),
\]
其中 \(\Omega\) 是一个固定的环形截断函数，在
\(W_1,W_2\) 的支集之积上等于一，且
\[
 H_{D,N}(s)=\Omega(s)
   [1-V_{\le}(2DNs/D_*)]e^{-DNs/Y_*}
   V_{\rm term}(DNs/U^{21}).
\]
每个固定的对数导数在
\(D,N,U\) 中一致有界。在截断过渡处其自变量有界，而
指数的对数导数等于其自变量的一个多项式乘以该指数。
对数 Fourier 反演因此给出
\[
 \begin{split}
 H_{D,N}(xy)&=\frac1{2\pi}\int_{\mathbb R}
      \widehat H_{D,N}(\nu)x^{i\nu}y^{i\nu}\,d\nu,\\
 \int_{|\nu|>cT_1}|\widehat H_{D,N}(\nu)|
       (1+|\nu|)^J\,d\nu&\ll_{J,N_0,c}T_1^{-N_0}
 \end{split}
\]
对每个固定的 \(J,N_0\)。整条直线上的系数 \(L^1\)-范数
也一致有界。两个
有限多项式的平凡界具有固定的 \(U\)-幂，因为它们的长度
被 \(U^{21+o(1)}\) 界住。在 \(\tau>0\) 固定之后，选取
\(N_0\) 使得被丢弃的 Fourier 尾部为 \(o(1)\)。

剩余积分之和的绝对值有
正的下界。二进对的数目与
一致的 \(L^1\)-范数表明，对一对以及一个
\(|\nu|\le cT_1\)，两个未正规化块之积具有
绝对值 \(\gg(\log U)^{-2}\)。两个剖面都具有
相同的高度 \(\gamma-\nu\)：
\[
 \begin{split}
 W_M(x)&=W_1(x)V_{\le}(Dx/D_*)x^{-\sigma-i(\gamma-\nu)},\\
 W_S(y)&=W_2(y)y^{-\sigma-i(\gamma-\nu)}.
 \end{split}
\]
第一个额外的截断函数仍属于逆环形剖面
且不产生第二个频率。通过中心正规化移除因子
\(D^{1/2-\sigma}\) 与 \(N^{1/2-\sigma}\)
得到
\[
 |M_rS_m|^2\gg(\log U)^{-4}(DN)^{2\sigma-1}
 \ge U^{\delta(r+m)-\epsilon}
\]
对充分大的 \(U\)。支集不等式给出
方程~\eqref{eq:detector-note} 中的两个长度界。

对这些剖面应用引理~\ref{lem:detector-dyads}，选取其
损失小于当前的 \(\epsilon\)。将
乘积下界与两个上界比较给出
\[
 \delta\{m-\min(m,1-m)\}
   =2\delta(m-1/2)_+\ll\epsilon.
\]
由于 \(\delta\ge1/50\)，这蕴含
\(m\le1/2+O(\epsilon)\)，其中常数为绝对常数。
支集不等式随后给出
\(r\ge t-1/2-O(\epsilon)\)。
将乘积下界依次除以各个上
界，即得
方程~\eqref{eq:saturated-witness} 中的两个各自下界。只有二进支集
误差是 \(O(1/\log U)\)；所有其他损失都是任意小的
固定幂。引理~\ref{lem:smooth-calculus} 允许所述的
逐行选择，且带有一个固定的多项式高度因子。
\end{proof}

\section{六次筛法的行计数}\label{sec:sextic-row-count}

零点探测器为底箱之上的每一行提供一个大的逆多项式。我们现在把该下界
转化为物理行的计数。我们首先按主约定证明所需的无平方因子六次大筛法，
使用 Blomer、Goldmakher 与 Louvel~\cite[第 3 节]{BGL} 的范数递归方法。然后我们
固定一个物理行中素因子赋值至少为二的部分，并把该筛法应用于其无平方因子
因子。固定部分的大小也给出行数的一个独立上界；两个界中较好的一个给出所需的
指数。

\subsection{主约定下的六次大筛法}

我们继续使用固定的有限素理想集合 $S$，它包含
$6$ 之上的素理想以及所有固定的导子素理想。下一个引理中所有理想指标
都与 $S$ 互素，并由其主生成元表示。特别地，“无平方因子”指的是
$\mathcal O=\mathbb Z[\omega]$ 的理想，而非其有理范数。

\begin{lemma}[六次大筛法]\label{lem:sextic-large-sieve}
设 $K,D\ge1$，设 $\varsigma\in\{-1,1\}$，并设 $(c_n)$ 为以 $S$ 之外的无平方因子理想为指标的任意
复数列。该数列独立于行 $k$ 而固定。对每个
$\epsilon>0$，
\begin{equation}\label{eq:sextic-large-sieve}
 \sum_{\substack{k\ {\rm sf}\\q_k\le K}}
 \left|\sum_{\substack{n\ {\rm sf}\\q_n\le D}}
       c_n\chi_n(k)^{\varsigma}\right|^2
 \ll_{S,\epsilon}(KD)^\epsilon
       \bigl\{K+D+(KD)^{2/3}\bigr\}
       \sum_{\substack{n\ {\rm sf}\\q_n\le D}}|c_n|^2.
\end{equation}
本公式中的每个符号在非单位元上取原来的零值。
允许对行集合作固定的限制。允许通过对该系数列作零延拓来
限制系数支集。这样的限制可以依赖于在行和之前固定的对象，
但系数不得在其他方面依赖于 $k$。
\end{lemma}

\begin{proof}
我们遵循 Blomer、Goldmakher 与
Louvel \cite[第~3 节]{BGL} 的高阶大筛法递归，为当前的
主正规化作局部验证。在有限 Fourier 计算中，
各个剩余特征作用于元素；它们的单位平凡商
就是我们把 Poisson 求和应用于其上的理想特征。

对由主生成元表示的好理想 $a,b$，令
\[
 F_b(a)=\chi_a(b),\qquad J=\{1,2,4\}.
\]
这里使用的是原来的零延拓符号，即便当 $a$ 非无平方因子时也是如此。
它在 $a$ 与 $b$ 中都是乘性的。在下面的每个幂中，包括六的倍数，
非单位元仍然取值为零。
对 $M>0$ 令
\[
 \mathcal I_M=\{a\text{ good}:M<q_a\le2M\},\qquad
 \mathcal A_M=\{a\in\mathcal I_M:a\text{ squarefree}\}.
\]
对支集在 $\mathcal A_N$ 上的列写
$\|\lambda\|_2^2=\sum_{b\in\mathcal A_N}|\lambda_b|^2$，并定义
\begin{align*}
 D_j(M,N)&=\sup_{\|\lambda\|_2=1}
  \sum_{a\in\mathcal A_M}
  \left|\sum_{b\in\mathcal A_N}\lambda_bF_b(a)^j\right|^2,\\
 E_j(M,N)&=\sup_{\|\lambda\|_2=1}
  \sum_{a\in\mathcal I_M}
  \left|\sum_{b\in\mathcal A_N}\lambda_bF_b(a)^j\right|^2.
\end{align*}
空支集给出范数零。这些范数包含所有有限射线
类。我们将为 $D_1$ 证明对称界；引理中的矩阵是它的转置，
且两个长度互换。

无平方因子范数 $D_j$ 是我们的目标。我们还需要 $E_j$，因为
Poisson 求和会产生任意的求值理想。在其定义和中展开平方
并对列特征配对，将把行长 $M$ 替换为一个阶为 $N^2/M$ 的对偶
长度。然后幂分解把无限制的行和重新化为无平方因子范数。
它选取一次幂或二次幂因子，因此指数 $j\in\{1,2,4\}$ 被保持在一起：
该集合在 $j\mapsto2j\pmod6$ 下封闭。我们将利用这一循环来
改进 $D_j$ 界中的一个临时指数。

\medskip\noindent\emph{配对求和公式。}
使用固定的有限群
\[
 \mathcal T=\ker\bigl((\mathcal O/36\mathcal O)^\times
              \longrightarrow(\mathcal O/3\mathcal O)^\times\bigr),
 \qquad t(b)=b\bmod36\mathcal O.
\]
主正规化将它们与模
$36\mathcal O$ 的射线类等同起来：每个单位轨道都有一个模 $3$ 等于一的
唯一代表。特别地 $t$ 是乘性的。这一细分
是固定的，因为 $2$ 与 $3$ 之上的素理想属于 $S$。

对全局单位 $u$ 与好素理想 $p$，模 $p$ 的约化给出
$\chi_p(u)=u^{(q_p-1)/6}$。由于 $q_p\equiv1\pmod6$，对
$b$ 的素因子相乘给出
\begin{equation}\label{eq:sextic-sieve-unit-class}
 \chi_b(u)=u^{(q_b-1)/6}.
\end{equation}
指数按模六来读：展开整数
$1+6h$ 的乘积即证明指数的相等。因此相等的 $t$-类在
六个全局单位上具有相同的限制。它们也具有相同的模
$4$ 剩余，这控制着引理~\ref{lem:fixed-gauss-phase} 的互反因子
$\mathcal R$。

设 $b,c$ 为互素的无平方因子好理想，满足 $t(b)=t(c)$，并设
$j\in J$。元素上的原始特征
\[
 \Theta_{b,c}^{(j)}(z)=\chi_b(z)^j\overline{\chi_c(z)^j}
       \quad (z\bmod bc)
\]
由方程~\eqref{eq:sextic-sieve-unit-class} 在全局单位上平凡。
因此它通过
\[
 \Psi_{b,c}^{(j)}((z))=\Theta_{b,c}^{(j)}(z).
\]
在理想上定义一个有限阶射线特征。
对与 $bc$ 互素的分数理想，剩余符号在局部单位分数上求值；
同样的单位相消使该定义与生成元无关。在每个好素理想处，六次剩余
特征具有恰好为六的阶，因为剩余乘法群是
循环的。因此其幂 $j$ 与 $-j$ 非平凡。CRT 表明
$\Theta_{b,c}^{(j)}$ 具有原始导子 $bc$，且该射线特征具有
相同的导子：略去一个素理想将会仅变化其剩余从而与这一非平凡性矛盾。
它在与 $bc$ 相交的理想上作零延拓。
对 $b=c=1$ 的一对，定义 $\Psi_{1,1}^{(j)}$ 为主特征。

六次互反律，其因子由 $t(b)=t(c)$ 固定，对每个
好理想 $a$ 给出
\begin{equation}\label{eq:sextic-sieve-pair-values}
 \Psi_{b,c}^{(j)}(a)=F_b(a)^j\overline{F_c(a)^j}.
\end{equation}
对互素的输入两个互反因子相消；在其他输入上
两边均为零。利用算术预备中已定义的 Gauss 和，使用 $z=cv+bw$ 的 CRT 给出
\begin{equation}\label{eq:sextic-sieve-pair-gauss}
 \begin{split}
 \Gamma_{b,c}^{(j)}
 &:=q_{bc}^{-1/2}\sum_{z\bmod bc}\Theta_{b,c}^{(j)}(z)e(z/(bc))\\
 &=\chi_b(c)^j\overline{\chi_c(b)^j}\gamma_j(b)\gamma_{-j}(c)
  =\eta_{t(b),j}\gamma_j(b)\gamma_{-j}(c),
 \end{split}
\end{equation}
其中 $\eta_{t(b),j}=\mathcal R(b,c)^j$ 只依赖于公共类。
这里每个 $\gamma_{\pm j}$ 由素 Gauss 正交性与 CRT 具有模一。
公式~\eqref{eq:sextic-sieve-pair-gauss} 保留了
$\gamma_{-j}$ 本身，包括它在共轭下 $-1$ 处的因子。

固定选取一个非负的 $W\in C_c^\infty((0,\infty))$，满足在
$[1,2]$ 上 $W\ge1$。对算术预备中的 Fourier 变换与自对偶测度，
令 $w(z)=W(|z|^2)$ 并通过
\[
 \widehat w(y)=\mathscr W(|y|^2).
\]
定义 $\mathscr W$。
该变换是径向的、在零处光滑且快速递减。于是
\begin{equation}\label{eq:sextic-sieve-transform-bounds}
 \mathscr W(x)\ll_A(1+x)^{-A},\qquad
 \int_{\mathbb R}|\mathcal M\mathscr W(\sigma+it)|\,dt<\infty
       \quad(\sigma>0).
\end{equation}
这里 $\mathcal M$ 表示 Mellin 变换。第二个断言
由 Mellin 积分中的分部积分得到；所有对数
导数在零处有界、在无穷处快速
递减。对 Fourier 变换作伸缩，$z\mapsto W(|z|^2/M)$ 给出
$M\mathscr W(M|y|^2)$。

对 $bc\ne1$，本原原始特征的有限 Fourier 反演为
\[
 \sum_{v\bmod bc}\Theta_{b,c}^{(j)}(v)e(vy/(bc))
 =\overline{\Theta_{b,c}^{(j)}(y)}\sqrt{q_{bc}}\,
       \Gamma_{b,c}^{(j)}.
\]
对非单位元 $y$，左边为零，因为一个非平凡的局部
特征被对平凡的加性特征求和。对所给的
$e$ 与 $d\mu$，基 $(1,\omega)$ 上的迹配对具有矩阵
$\left(\begin{smallmatrix}0&1\\1&-1\end{smallmatrix}\right)$，行列式为
$-1$，且 $\mathcal O$ 的 $d\mu$-余体积为一。于是
$bc\mathcal O$ 的对偶为 $(bc)^{-1}\mathcal O$，而
$bc\mathcal O$ 的余体积为 $q_{bc}$。
陪集 Poisson 求和，随后除以每个理想的六个生成元，
因此给出
\begin{equation}\label{eq:sextic-sieve-pair-poisson}
 \sum_{\mathfrak a\ne0}W(q_{\mathfrak a}/M)\Psi_{b,c}^{(j)}(\mathfrak a)
 =\frac{M\Gamma_{b,c}^{(j)}}{\sqrt{q_{bc}}}
   \sum_{\mathfrak a\ne0}\mathscr W(Mq_{\mathfrak a}/q_{bc})
                 \overline{\Psi_{b,c}^{(j)}(\mathfrak a)}.
\end{equation}
这里两个和都在所有非零整理想上取。除以六
利用了 $\Theta_{b,c}^{(j)}$ 的单位平凡性与权重的径向性。
零频率消失，因为该特征非平凡。

\medskip\noindent\emph{平方范数的 Poisson 变换。}
对支集在 $\mathcal A_N$ 中一个 $t$-类上的 $\lambda$，令
\[
 Q_j(M,N;\lambda)=
 \sum_{\substack{b,c\in\mathcal A_N\\(b,c)=1}}
   \lambda_b\overline{\lambda_c}
   \sum_{\mathfrak a\ne0}W(q_{\mathfrak a}/M)
                         \Psi_{b,c}^{(j)}(\mathfrak a).
\]
我们断言，对每个 $\varepsilon>0$，
\begin{equation}\label{eq:sextic-sieve-pair-norm}
 \begin{split}
 |Q_j(M,N;\lambda)|
 &\ll_{S,\varepsilon}(2+MN)^\varepsilon\|\lambda\|_2^2\\
 &\quad\times
 \left\{M+\frac MN
   \max_{\substack{1\le H\\
       H\ll_{S,\varepsilon}(2+MN)^\varepsilon N^2/M}}
       E_j(H,N)\right\}.
 \end{split}
\end{equation}
最大值在二进壳层上取，若为空则为零。若 $N\le1$，
只可能出现单位好理想，环形理想计数即证明
$M$ 项。此外当 $M$ 低于某个固定正常数时 $Q_j=0$，
因为 $W$ 具有固定的支集上端点，而非零理想的范数
至少为一。

设 $N>1$ 并正规化 $\|\lambda\|_2=1$。则 $Q_j$ 中的每一对
都是非平凡的且 $q_{bc}\asymp N^2$。应用
方程~\eqref{eq:sextic-sieve-pair-poisson} 并对对偶理想取三角
不等式。对小的 $\delta>0$ 令
$Z=(2+MN)^\delta$。Cauchy 给出 $(\sum_b|\lambda_b|)^2\ll N$。
由方程~\eqref{eq:sextic-sieve-transform-bounds}，
$q_{\mathfrak a}>ZN^2/M$ 的部分，对任意 $A>1$，至多为
\[
 M\sum_{q_{\mathfrak a}>ZN^2/M}(Mq_{\mathfrak a}/N^2)^{-A}
 \ll N^2Z^{1-A}.
\]
若截断低于一，同样的界成立。取 $A$ 大便使
其成为任意负幂误差。

在剩余的求和中写 $\mathfrak a=a_0\mathfrak s$，其中 $a_0$ 为好理想、
$\mathfrak s$ 的支集在 $S$ 上，并为每个固定的 $\mathfrak s$ 选取一个生成元 $s_0$。
方程~\eqref{eq:sextic-sieve-pair-values}
与~\eqref{eq:sextic-sieve-pair-gauss} 给出分离的因子
\[
 \overline{\Psi_{b,c}^{(j)}(a_0\mathfrak s)}
 =\overline{F_b(a_0)^j}F_c(a_0)^j
       \overline{\chi_b(s_0)^j}\chi_c(s_0)^j.
\]
所有 $\chi_b(s_0)$ 具有模一。将 $a_0$ 分成壳层
$a_0\in\mathcal I_H$，其中 $H\ll ZN^2/(Mq_{\mathfrak s})$，并
在 $\Re w=\delta$ 上插入 $\mathscr W$ 的
Mellin 反演。除一个依赖于
$M,q_{a_0},q_{\mathfrak s}$ 的有界因子外，两个
系数列为
\[
 A_b=\lambda_b\gamma_j(b)\overline{\chi_b(s_0)^j}
                    q_b^{-1/2+\delta+it},\qquad
 B_b=\overline{\lambda_b}\gamma_{-j}(b)\chi_b(s_0)^j
                    q_b^{-1/2+\delta+it}.
\]
它们在 $a_0$ 求和过程中固定，且具有绝对值
$\ll N^{-1/2+\delta}|\lambda_b|$。通过 M\"obius 反演
检测 $(b,c)=1$ 并在 $a_0$ 中应用 Cauchy 给出
\begin{align*}
 &\sum_{a_0\in\mathcal I_H}
  \left|\sum_{\substack{b,c\in\mathcal A_N\\(b,c)=1}}
       A_bB_c\overline{F_b(a_0)^j}F_c(a_0)^j\right|\\
 &\quad\le E_j(H,N)\sum_{\mathfrak d}
   \left(\sum_{\mathfrak d\mid b}|A_b|^2\right)^{1/2}
   \left(\sum_{\mathfrak d\mid b}|B_b|^2\right)^{1/2}
 \ll N^{-1+3\delta}E_j(H,N).
\end{align*}
这里除数界吸收了关于 $\mathfrak d$ 的求和。对整个
第一个多项式取共轭使用相同的 $E_j$ 范数。满足
$H<1$ 的因子只有单位求值理想，并满足 $E_j(H,N)\ll N$。
$\mathfrak s$ 与二进壳层的选择只有 $\log(2+ZN^2)$ 的固定幂次
多种，因为 $S$ 固定且非零
$Q_j$ 有 $M$ 有下界。方程~\eqref{eq:sextic-sieve-transform-bounds}
对 $t$ 变量求积分。以 $\varepsilon$ 为准取 $\delta$ 足够小
即证明方程~\eqref{eq:sextic-sieve-pair-norm}，
包括其 $M$ 项上允许的全局损失。

我们现在把这一配对估计应用于 $E_j$。按
$t$ 划分其内和，插入 $W$，并展开平方。将两个无平方因子指标写为
$db,dc$，其中 $d$ 是它们的 gcd、$(b,c)=1$ 且 $(d,bc)=1$。它们的剩余
类一致，因为 $t$ 是乘性的。指标乘性与
方程~\eqref{eq:sextic-sieve-pair-values} 在每个好理想 $a$ 上给出
\[
 F_{db}(a)^j\overline{F_{dc}(a)^j}
   =1_{(a,d)=1}\Psi_{b,c}^{(j)}(a).
\]
令 $\mathfrak s_S=\prod_{\mathfrak p\in S}\mathfrak p$。对 $(a,\mathfrak s_Sd)=1$ 作 M\"obius 反演
将 $a$ 写为 $a=\mathfrak r\mathfrak a$，其中
$\mathfrak r\mid\mathfrak s_Sd$。为每个固定的 $\mathfrak r$ 选取一个生成元 $r_0$。
$\Psi$ 在 $\mathfrak r$ 处的值分离为
$\chi_b(r_0)^j\overline{\chi_c(r_0)^j}$，于是公共剩余列为
\[
 \lambda'_b=\lambda_{db}1_{(b,d)=1}\chi_b(r_0)^j,
 \qquad |\lambda'_b|\le|\lambda_{db}|.
\]
它的支集在 $\mathcal A_{N/q_d}$ 中一个类上。以长度 $M/q_{\mathfrak r}$
与 $N/q_d$ 应用方程~\eqref{eq:sextic-sieve-pair-norm}。求和后的系数
质量为至多
\[
 \sum_d\sum_{\mathfrak r\mid\mathfrak s_Sd}\sum_b|\lambda_{db}|^2
 \le d_{\mathcal O}(\mathfrak s_S)
       \sum_m d_{\mathcal O}(m)^2|\lambda_m|^2
 \ll_{S,\delta}(2+N)^\delta\|\lambda\|_2^2.
\]
写 $\Delta_{\mathrm r}=q_{\mathfrak r}$ 与 $\Delta_{\mathrm d}=q_d$，我们有
$1\le\Delta_{\mathrm d}\le2N$ 与 $1\le\Delta_{\mathrm r}\le q_{\mathfrak s_S}\Delta_{\mathrm d}$。
我们得到递归接口
\begin{equation}\label{eq:sextic-sieve-norm-reflection}
 \begin{split}
 E_j(M,N)
 &\ll_{S,\varepsilon}(2+MN)^\varepsilon
 \max_{\substack{1\le\Delta_{\mathrm d}\le2N\\
                 1\le\Delta_{\mathrm r}\le q_{\mathfrak s_S}\Delta_{\mathrm d}}}
 \left\{\frac M{\Delta_{\mathrm r}}\right.\\
 &\qquad\left.+
 \frac{M\Delta_{\mathrm d}}{N\Delta_{\mathrm r}}
 \max_{\substack{1\le H\\
       H\ll_{S,\varepsilon}(2+MN)^\varepsilon
              N^2\Delta_{\mathrm r}/(M\Delta_{\mathrm d}^2)}}
       E_j(H,N/\Delta_{\mathrm d})\right\}.
 \end{split}
\end{equation}
最大值可以限制在非空范围上。这是 \cite[引理~3.2--3.3]{BGL} 中
gcd 与 Poisson 步骤的组合；这里使用的每个序列在其范数
和内都是公共的，且 $S$ 保持固定。

\medskip\noindent\emph{初始界与幂界。}
行包含给出 $D_j\le E_j$。Hilbert 空间对偶性与六次
互反律，按有限 $t$-类分裂一个矩阵变量，给出
\begin{equation}\label{eq:sextic-sieve-norm-symmetry}
 D_j(M,N)\ll_S D_j(N,M).
\end{equation}
相位取单位值，且非互素元素在两个方向上都为零。
我们还有初始估计
\begin{equation}\label{eq:sextic-sieve-norm-initial}
 D_j(M,N)\ll_S M^2+N\qquad(M,N\ge1).
\end{equation}
事实上，对无平方因子行 $a$，有限 Fourier 反演在每个 $b$ 上，
包括非单位元，给出
\[
 \chi_a(b)^j=\frac1{\sqrt{q_a}\gamma_{-j}(a)}
       \sum_{x\bmod a}\overline{\chi_a(x)^j}e(xb/a).
\]
Cauchy 至多耗费 $\#(\mathcal O/(a))^\times/q_a\le1$。来自
$a\in\mathcal A_M$ 的既约分数 $x/a$ 在模 $\mathcal O$ 下互不相同
且彼此间隔
至少 $1/(2M)$：$x/a-y/c-h$ 的一个非零分子具有绝对值
至少为一，而 $|ac|\le2M$；模 $\mathcal O$ 的相等迫使
相同的主分母与剩余。引理~\ref{lem:planar-additive-sieve} 的对偶，
其系数生成元位于球 $q_b\le2N$ 内，证明了方程~\eqref{eq:sextic-sieve-norm-initial}。

